\section{How the design was found}
\label{sec:history}

The design came from three build-and-evaluate cycles: two earlier configurations, each built and then tested in a pre-registered study, and the product, tested in the third, Study~3 (Section~\ref{sec:test}). Neither earlier configuration was better than Rules in context; after each study we attributed its shortfalls, post hoc, to likely causes and changed the design. The Supplement gives both studies in full (Sections~\ref{S-app:study1} and~\ref{S-app:study2}) and their deviations and robustness analyses (Sections~\ref{S-supp:deviations} and~\ref{S-supp:robustness}).

\paragraph{Study 1}
The Study~1 configuration made one reading, had a second model call accept or reject each value, kept only values whose quotations were in the file, and evaluated the workflow in strong Kleene logic and, as a variant on the same readings, by exact deferral; the tenancy figures below are those of exact deferral. It was tested on 90 tenancy claims and on 27 pairs of theft texts that differ in one fact, three runs each. The theft workflow had been extracted by a model from insurance law and policy terms. With the frontier model, the primary analysis found it no better than Rules in context: \thTerraCPWrong{} against \thTerraRWrong{} wrong theft changes ($p=\ptTerraHThreeRaw$) and \ctTerraHOneA{} against \ctTerraHOneB{} unsafe errors per tenancy claim ($p=\ctTerraHOneP$). A sensitivity analysis decided after the runs, which counts documents handed to a person as missed, found it worse on tenancy (\ctTerraHOneStrictA{} against \ctTerraHOneStrictB, $p=\ctTerraHOneStrictP$). With the cost-efficient model it made fewer unsafe errors per tenancy claim (\ctLunaHOneA{} against \ctLunaHOneB, $p=\ctLunaHOneP$) and fewer wrong theft changes, \thLunaCPWrong{} against \thLunaRWrong; most of the comparator's theft errors were one error repeated across a story, and counted once they were \ctLunaHThreeDedupRA{} against \ctLunaHThreeDedupRB{} per pair ($p=\ctLunaHThreeDedupRP$). Because code applied the rules, each error could be traced to a rule or a reading (Table~\ref{tab:history}).

\begin{table}[t]
\caption{Likely causes found in Study~1, the change each led to, and the evidence for the change. Counts are from Study~1 with the frontier model unless stated; Study~2 values are with the cost-efficient model. ``Recomputed'' marks evidence computed again from recorded readings, without new model calls. No row's evidence comes from a confirmatory test.}
\label{tab:history}
\small
\begin{tabularx}{\linewidth}{@{}>{\raggedright\arraybackslash}X>{\raggedright\arraybackslash}X>{\raggedright\arraybackslash}X@{}}
\toprule
Cause in Study 1 & Change & Evidence \\
\midrule
Faulty rules: \ecWrongFaulty{} of \thTerraCPWrong{} wrong theft changes came from \rvImplicated{} of the \rvRules{} extracted rules. & A review of every rule against the passages it cites, after a repair without labels failed. & Study~2, new runs: \stThCPvTwoLunaErrors{} errors per theft pair with the label-free repair; recomputed: \stAblThOriginalPackErrors{} with the original rules ($p=\stAblThOriginalPackVsP$). The review itself was evaluated on development data only. \\
Silence read as doubt: \jaUnresFalse{} of \jaUnres{} unresolved values concerned facts that were false in the scenario, which the claims leave unmentioned. & Presumptions with a source clause. & Study~2, recomputed: \stTenCPvTwoLunaHeldReq{} held needed documents per claim with presumptions, \stAblTenNoPresumptionHeldReq{} without. \\
The checking call rejected \verRejCondTrue{} of \verCheckedCondTrue{} true values, all correct, and none of the \jaFlipped{} wrong ones. & A second reading that must agree; no checking call. & Study~2, recomputed by pairing each run's reading with another run's reading of the same text: errors per theft pair fell from \stAblThTwoReadingsVsB{} to \stAblThTwoReadingsVsA{} ($p=\stAblThTwoReadingsVsP$, the smallest attainable). \\
Strong Kleene evaluation, the production engine's, held \tenTerraCPkHeld\% of needed tenancy documents. & Exact deferral, which Study~1 had evaluated as a variant, became the evaluation. & Study~1, on the same readings: exact deferral held \tenTerraExactCut\% (frontier) and \tenLunaExactCut\% (cost-efficient) fewer, with no unnecessary request added. \\
\bottomrule
\end{tabularx}
\end{table}

\paragraph{Study 2}
Study~2 confirmed none of its four primary contrasts. It tested the second configuration on sealed new data, 72 tenancy claims and 36 theft pairs. The configuration made one reading without a checking call, used presumptions and exact deferral, and executed a theft workflow repaired without labels. Against Rules in context with the same model and workflow, theft errors per pair were \stAA{} against \stAB{} (raw $p=\stAP$, Holm $p=\stAHolm$); tenancy showed no significant difference (\stBA{} against \stBB{} unsafe errors per claim). With the cost-efficient model it was worse on theft than the frontier model applying the original workflow in context, \stCA{} against \stCB{} errors per pair ($p=\stCP$), and its non-inferiority to the frontier model on tenancy was not shown ($p=\stDPNI$).

\paragraph{What failed}
A first redesign withheld every rule that an automatic check could not repair; it was dropped on development data before Study~2. Repair without labels also failed in Study~2: a model checked each extracted rule against its cited passages and repaired flagged rules with typed operations. The repaired workflow did not lower the configuration's errors (Table~\ref{tab:history}), and it raised those of Rules in context with the frontier model from \stThRTerraErrors{} to \stThRVTwoTerraErrors{} per pair. Learning from the same handler corrections did not separate the configuration from Rules in context with the corrections as memory (\stEvoAfterA{} against \stEvoAfterB{} errors per theft pair, $p=\stEvoAfterP$).

\paragraph{What remains}
Four changes from Study~1 remain: a reviewed workflow, readings that must agree, presumptions and exact deferral. Together they form the decision rule of Section~\ref{sec:rule}. The review checked each of the \rvRules{} extracted theft rules against the passages it cites. Two AI editor sessions worked independently, and a change was adopted only when both made it, or one made it and the other approved it; 9 rules changed (one was split, giving 99), 14 were queued for a jurist, and no jurist has reviewed them. The editors knew the theft conditions that Study~3 reuses. The joint effect of the four changes was measured only on development data before Study~3.

\paragraph{A public benchmark}
An unplanned pilot on ShARC \cite{saeidi2018sharc} after Study~2 compared the two-reading core with prompting on short single-rule items; no claim in this article rests on it (Supplement, Section~\ref{S-supp:sharc}).
